\documentclass[11pt]{article}
\usepackage[utf8]{inputenc}
\usepackage{amsmath,amssymb}
\usepackage{hyperref}
\title{Robust Single Cell Rna Sequencing Analysis under Distribution Shift}
\author{Anonymous}
\date{}
\begin{document}
\maketitle

\begin{abstract}
This paper studies Single Cell Rna Sequencing Analysis. Grounded in a real, citable literature \cite{ref1,ref2,ref3,ref4,ref5,ref6,ref7,ref8},
we propose a focused method and report \textbf{illustrative} experiments.
All references below are real and verifiable (OpenAlex/arXiv); the reported
numbers are simulated to illustrate intended behaviour.
\end{abstract}

\section{Introduction}
Single Cell Rna Sequencing Analysis is an active research area. This work targets a concrete gap identified from
the recent literature and contributes a method together with an evaluation protocol.

\section{Related Work (real literature)}
The following works, retrieved from public bibliographic APIs, frame our study:
\begin{itemize}
\item Prior work (1990) — "PCR protocols — A guide to methods and applications", Trends in Genetics (cited 11865).
\item Jin et al. (2021) — "Inference and analysis of cell-cell communication using CellChat", Nature Communications (cited 8535).
\item Becht et al. (2018) — "Dimensionality reduction for visualizing single-cell data using UMAP", Nature Biotechnology (cited 5822).
\item Tang et al. (2019) — "GEPIA2: an enhanced web server for large-scale expression profiling and interactive analysis", Nucleic Acids Research (cited 5517).
\item Aran et al. (2019) — "Reference-based analysis of lung single-cell sequencing reveals a transitional profibrotic macrophage", Nature Immunology (cited 5430).
\item Wang et al. (2008) — "Alternative isoform regulation in human tissue transcriptomes", Nature (cited 5390).
\end{itemize}

\section{Method}
We model distribution shift explicitly and optimise a worst-case objective over an uncertainty set of plausible test distributions. We instantiate this idea for Single Cell Rna Sequencing Analysis and analyse its cost/quality trade-off.

\section{Experiments (Illustrative)}
\begin{center}
\begin{tabular}{lcc}
System & Primary Metric & Efficiency \\ \hline
Strong baseline & 0.812 & 1.00$\times$ \\
Ablation & 0.828 & 0.92$\times$ \\
\textbf{Ours} & \textbf{0.851} & \textbf{0.71$\times$}
\end{tabular}
\end{center}
\emph{Metrics simulated to illustrate the intended effect; not measured.}

\section{Conclusion}
We connected a real literature to a concrete method for Single Cell Rna Sequencing Analysis. Future work will
validate the illustrative results with real experiments.

\begin{thebibliography}{9}
\bibitem{ref1} . PCR protocols — A guide to methods and applications. \emph{Trends in Genetics}, 1990. DOI:10.1016/0168-9525(90)90186-a
\bibitem{ref2} Suoqin Jin, Christian F. Guerrero‐Juarez, Lihua Zhang et al.. Inference and analysis of cell-cell communication using CellChat. \emph{Nature Communications}, 2021. DOI:10.1038/s41467-021-21246-9
\bibitem{ref3} Étienne Becht, Leland McInnes, John Healy et al.. Dimensionality reduction for visualizing single-cell data using UMAP. \emph{Nature Biotechnology}, 2018. DOI:10.1038/nbt.4314
\bibitem{ref4} Zefang Tang, Boxi Kang, Chenwei Li et al.. GEPIA2: an enhanced web server for large-scale expression profiling and interactive analysis. \emph{Nucleic Acids Research}, 2019. DOI:10.1093/nar/gkz430
\bibitem{ref5} Dvir Aran, Agnieszka Looney, Leqian Liu et al.. Reference-based analysis of lung single-cell sequencing reveals a transitional profibrotic macrophage. \emph{Nature Immunology}, 2019. DOI:10.1038/s41590-018-0276-y
\bibitem{ref6} Eric T. Wang, Rickard Sandberg, Shujun Luo et al.. Alternative isoform regulation in human tissue transcriptomes. \emph{Nature}, 2008. DOI:10.1038/nature07509
\bibitem{ref7} Itay Tirosh, Benjamin Izar, Sanjay M. Prakadan et al.. Dissecting the multicellular ecosystem of metastatic melanoma by single-cell RNA-seq. \emph{Science}, 2016. DOI:10.1126/science.aad0501
\bibitem{ref8} Yuhan Hao, Tim Stuart, Madeline H. Kowalski et al.. Dictionary learning for integrative, multimodal and scalable single-cell analysis. \emph{Nature Biotechnology}, 2023. DOI:10.1038/s41587-023-01767-y
\end{thebibliography}
\end{document}
